\documentclass{article}
\usepackage{hyperref}
\begin{document}
\section*{Early full record versus late scoped reference}
\paragraph{Matched control.}
Run the original pinned SQLite/RWARE and SubT replay of
\href{https://github.com/eternal-november/eos-research-programmes/pull/814}{the radio handoff}
with three presentation policies. Each receives the same independently
delivered report, source closure and activation cut. The donor's worker
has actually died following durable commit; only a read-fenced original
store provides the reconstructed prefix and settled or pending record.
The early arm delivers the full status just after the motion prefix,
at the fixture's early link at time $-3$. It stores bytes without
authorizing action, then reopens the recipient database and installs
the same held key only at the original decision cut. The late arms
receive at time zero. The final four-byte outcome is decoded and
compared byte for byte; the physical step count is checked separately.
No new movement on a settled key is accepted. Absent and stale
statuses cannot authorize the old duty in any arm.

\paragraph{Exact eligible cases.}
For each 864-case arm, the late full record passes 130, the late
one-byte reference passes 150, and early full staging passes 180.
For this fixture the request is feasible exactly when the active
transfer debit plus four is within Q's limit, and the reply is
feasible exactly when the prior independent debit plus four is within
B's limit, conditional on valid authority, physical feasibility or a
cached result, and fixed native links. The strict one-second expiry
is unchanged. Early preparation moves the first debit outside the
late window. Compact preparation changes the size of the live debit
without changing its expiry. It does not dominate early full staging
on count; early full does not dominate compact on total radio bytes.
Do not equate counts in this authored grid with failure probabilities.

\paragraph{Scope of the one-byte mapping.}
With a preinstalled basis catalog, ordinal zero is pending and the
positive ordinal references the exact decoded row. The basis code
binds the already reconstructed complete actor sequence in the two
tested prefixes; it is not a general-purpose checksum or signature.
An actual alternate native event order reaches the same agent
positions but changes row two and is rejected. If the sealed record is
not available at the prefix cut, its route is not concurrently usable,
or activation cannot be held fixed, early staging is not an admissible
comparison. This replay confirms those controls only for the authored
native fixture. It does not establish such a link, early availability,
or authority in a new field history. No independent I1/I2 certification.
\end{document}
